\documentclass[12pt]{article}
\usepackage[utf8]{inputenc}
\usepackage[T1]{fontenc}
\usepackage{hyperref}
\title{Rufiji Language Resources for K-12 Curriculum Development}
\author{Compiled by Claude Sonnet 5.5 (claude-sonnet-5-5) for Andr\'as Kornai}
\date{}
\begin{document}
\maketitle

\section{Introduction}
The code rui (Glottocode rufi1234) was assigned to a supposed language called Rufiji in Tanzania. The coordinator's brief says the Ethnologue editors found no evidence for such a language and are pursuing its removal from ISO 639-3. My searches could not confirm or refute that claim directly. They did find several catalogues that still list the code. Wikipedia's Ndengereko article lists rui as a duplicate code beside ndg (Ndengereko), citing Ethnologue 2015. Rufiji is also a river and district in Tanzania, so the name may refer to speakers of neighbouring languages such as Ndengereko. No educational material specific to a Rufiji language was found, and none is cited here.

\section{Primary Educational Resources}
None found. No primer, textbook or school programme for a language called Rufiji turned up in search.

\section{Linguistic and Reference Materials}
These entries record the status of the code and do not describe a language.
\begin{itemize}
\item ISO 639:r (Wikipedia): \url{https://en.wikipedia.org/wiki/ISO_639:r}
\item Ndengereko language (Wikipedia), which notes rui as a duplicate code: \url{https://en.wikipedia.org/wiki/Ndengereko_language}
\item Glottolog entry rufi1234: \url{https://glottolog.org/resource/languoid/id/rufi1234.bigmap.html}
\item ASJP wordlist entry, an old compilation: \url{https://asjp.clld.org/languages/RUFIJI}
\end{itemize}

\section{Religious and Community Resources}
None found. Joshua Project lists rui as a language, but I did not verify what it contains: \url{https://joshuaproject.net/languages/rui}

\section{Digital and Audio Resources}
Global Recordings Network lists a ``Rufiji language'' page. I did not verify its contents, and it may be an artifact of the code: \url{https://globalrecordings.net/en/language/rui}

\section{Curriculum Development Context}
No curriculum can responsibly be built for this entry. I recommend checking the current ISO 639-3 change-request history and the Ethnologue entry. If rui is retired, the relevant Tanzanian languages (for example Ndengereko, ndg) should be surveyed under their own codes.
\end{document}